% Academic Curriculum Vitae
% Updated July 2026
% Lu Niu

\documentclass[11pt,a4paper]{article}

\usepackage[dvipsnames]{xcolor}
\usepackage[version=4]{mhchem}
\usepackage{fancyhdr}
\usepackage{geometry}
\usepackage{tabularx}
\usepackage{array}
\usepackage{enumitem}
\usepackage{crimson}
\usepackage{helvet}
\usepackage[hidelinks]{hyperref}
\usepackage{microtype}

\geometry{left=1.8cm,right=1.8cm,top=1.85cm,bottom=1.65cm}
\setlength{\parindent}{0pt}
\setlength{\parskip}{0pt}
\setlength{\headheight}{18pt}
\linespread{0.96}
\urlstyle{same}

\definecolor{shadecolor}{rgb}{0.0,0.5,0.8}
\definecolor{subtitlecolor}{rgb}{0.5,0.8,0.9}

\newcommand{\namefont}[1]{{\normalfont\bfseries\Huge #1}}
\newcommand{\myname}{Lu Niu}
\newcommand{\mydegree}{PhD Candidate in Computational Condensed Matter Physics}
\newcommand{\cvupdated}{July 2026}

\newcommand{\cvsection}[1]{%
    \vspace{0.32cm}%
    \noindent\colorbox{subtitlecolor}{%
        \parbox{\dimexpr\textwidth-2\fboxsep\relax}{%
            \textcolor{white}{\sffamily\large\quad #1}%
        }%
    }%
    \par\vspace{0.14cm}%
}

\newcommand{\cvpublication}[1]{%
    \hangindent=1.25em\hangafter=1 #1\par\vspace{0.10cm}%
}

\newcommand{\cvproject}[3]{%
    \textbf{#1}\par
    \hfill\textit{#2}\par
    #3\par\vspace{0.16cm}%
}

\pagestyle{fancy}
\fancyhf{}
\fancyhead[C]{%
    \colorbox{shadecolor}{%
        \parbox{\dimexpr\textwidth-2\fboxsep\relax}{%
            \textcolor{white}{\sffamily\large\quad CURRICULUM VITAE}%
        }%
    }%
}
\fancyfoot[C]{\thepage}
\renewcommand{\headrulewidth}{0pt}

\begin{document}

\begin{minipage}[t]{0.47\textwidth}
    \vspace{0pt}
    \centering
    \namefont{\myname}
\end{minipage}\hfill
\begin{minipage}[t]{0.49\textwidth}
    \vspace{0.12cm}
    \raggedright
    \emph{\mydegree}\\[0.06cm]
    \emph{Thesis submitted June 2026}
\end{minipage}

\vspace{0.28cm}

\begin{minipage}[t]{0.45\textwidth}
    \vspace{0em}
    Condensed Matter Theory Group\\
    School of Physics\\
    The University of Sydney\\
    Sydney, NSW 2006, Australia
\end{minipage}\hfill
\begin{minipage}[t]{0.51\textwidth}
    \vspace{0em}
    \begin{tabular}{@{}>{\bfseries}p{1.45cm}l@{}}
        Email: & \href{mailto:LukeNiu@outlook.com}{LukeNiu@outlook.com}\\
        Web: & \href{https://condensed.physics.sydney.edu.au/luke.html}{condensed.physics.sydney.edu.au/luke.html}\\
        GitHub: & \href{https://github.com/Photonico}{github.com/Photonico}\\
        Phone: & +61 451 116 402
    \end{tabular}
\end{minipage}

\cvsection{Research Interests}
Computational condensed matter physics; topological matter; physics of complex systems; and AI for physics.

\cvsection{Education}

\textbf{The University of Sydney}\hfill\textit{2021--2026}\par
\quad\emph{Doctor of Philosophy in Physics}\par
\begin{tabularx}{\textwidth}{@{}>{\bfseries}p{10em}X@{}}
Field: & Condensed Matter Physics\\
Thesis: & \emph{Exploring Quantum Properties of Light-Element Xenes: DFT Studies of the Electronic Structure and Optical Response for Boron and Beryllium Allotropes}\\
Primary supervisor: & Professor Catherine Stampfl\\
Associate supervisors: & Dr Carla Verdi
\end{tabularx}

\vspace{1em}
\textbf{University of Science and Technology Beijing}\hfill\textit{2015--2017}\par
\quad\emph{Master of Research in Physics}\par
\begin{tabularx}{\textwidth}{@{}>{\bfseries}p{10em}X@{}}
Field: & Atomic and Molecular Physics\\
Thesis: & \emph{Effect of External Field on the I--V Characteristics through the Molecular Nano-junction}\\
Supervisor: & Professor Luxia Wang
\end{tabularx}

\cvsection{Expertise and Technical Strengths}
\begin{tabularx}{\textwidth}{@{}>{\bfseries}p{4.0cm}X@{}}
Theory and methods: & Density functional theory; quantum mechanics; many-body physics; linear response; Migdal--Eliashberg theory; phonons; topological band analysis\\
Materials modelling: & Atomic-structure optimisation; electronic bands and densities of states; dielectric and optical properties; magnetism; superconductivity; two-dimensional materials and heterostructures\\
Programming: & Python; \LaTeX\\
Software and computing: & VASP; Linux; Git; PyTorch; high-performance computing, including NCI Australia\\
Languages: & English; Mandarin Chinese (native)
\end{tabularx}

\cvsection{Teaching and Tutoring Experience}

\textbf{The University of Sydney}\par
Tutor for a graduate-level course in quantum mechanics and mathematical methods during the PhD candidature. Developed detailed teaching notes covering linear algebra, operator formalism, and mathematical methods used in quantum mechanics.

\vspace{0.12cm}
\textbf{University of Science and Technology Beijing}\par
\begin{tabularx}{\textwidth}{@{}Xr@{}}
Teaching Assistant in Analytical Mechanics & \textit{Spring 2016; Spring 2017}\\
Teaching Assistant in College Physics & \textit{Autumn 2015; Autumn 2016}
\end{tabularx}

\newpage
\cvsection{Selected Research Experience}

\cvproject{Electronic and optical properties of graphene/boron-based two-dimensional heterostructures}{Published 2024}{First-principles investigation of Graphene--\ce{BC3}, Graphene--Borophene, and Graphene--\ce{B4C3}, including structural properties, band structures, densities of states, interfacial charge redistribution, Schottky barriers, dielectric response, and derived optical properties.}

\cvproject{Functional properties of penta-bipyramid boron phases}{Published 2026}{First-principles study of bulk $o$-\ce{B14}, monolayer $o$-\ce{B14}, bilayer $o$-\ce{B14}, and the hydrogen-terminated bilayer. The work addresses semiconductivity, monolayer magnetism, superconductivity in the hydrogen-terminated bilayer, and anisotropic optical and plasmonic response.}

\cvproject{Topological and optical properties of pristine and hydrogen-functionalized beryllene}{Manuscript in preparation}{Density functional theory study of pristine $\alpha$- and $\beta$-beryllene, a cubic trilayer phase, and hydrogen-functionalized derivatives, covering structural and dynamical stability, electronic properties, parity-based $Z_{2}$ topology, dielectric response, and derived optical properties.}

\cvsection{Publications}

\textbf{2026}\par
\cvpublication{\textbf{Lu Niu}, Oliver J. Conquest, Hongyang Ma, Carla Verdi, and Catherine Stampfl, ``Superconductivity, Magnetism, and Directional Plasmonic Modes in Two-Dimensional Penta-Bipyramid Boron Allotropes,'' \emph{Physical Review Materials} \textbf{10}, 074001 (2026). \href{https://doi.org/10.1103/b2q9-fytt}{doi:10.1103/b2q9-fytt}.}

\textbf{2024}\par
\cvpublication{\textbf{Lu Niu}, Oliver J. Conquest, Carla Verdi, and Catherine Stampfl, ``Electronic and Optical Properties of 2D Heterostructure Bilayers of Graphene, Borophene and 2D Boron Carbides from First Principles,'' \emph{Nanomaterials} \textbf{14}, 1659 (2024). \href{https://doi.org/10.3390/nano14201659}{doi:10.3390/nano14201659}.}

\textbf{2018}\par
\cvpublication{\textbf{Lu Niu} and Luxia Wang, ``Effect of External Field on the I--V Characteristics through the Molecular Nano-junction'' (in Chinese), \emph{Acta Physica Sinica} \textbf{67}, 027304 (2018). \href{https://doi.org/10.7498/aps.67.20171604}{doi:10.7498/aps.67.20171604}.}

\textbf{2017}\par
\cvpublication{Dandan Zhao, \textbf{Lu Niu}, and Luxia Wang, ``Plasmon Enhanced Heterogeneous Electron Transfer with Continuous Band Energy Model,'' \emph{Chemical Physics} \textbf{493}, 194--199 (2017). \href{https://doi.org/10.1016/j.chemphys.2017.07.005}{doi:10.1016/j.chemphys.2017.07.005}.}

\cvsection{Manuscript in Preparation}
\cvpublication{\textbf{Lu Niu}, Oliver J. Conquest, Hongyang Ma, Carla Verdi, and Catherine Stampfl, ``Emergent Topological and Optical Properties in Pristine and Hydrogen-Functionalized Two-Dimensional Beryllium Phases,'' manuscript in preparation (2026).}

\cvsection{Conferences and Academic Activities}

\begin{tabularx}{\textwidth}{@{}>{\bfseries}p{1.25cm}X@{}}
2025 & The 47th Annual Condensed Matter and Materials Meeting (``Wagga 2025''), 4--7 February 2025, The University of Queensland, Brisbane, Australia.\\[0.08cm]
2023 & NSW Student Computational Chemistry Meeting, 10--12 July 2023, The Sebel Harbourside Kiama, Kiama, Australia.\\[0.08cm]
2016 & The 2nd Joint Workshop on Condensed Matter Science, Peking University and IMPRS, 14--18 November 2016, Peking University, Beijing, China.
\end{tabularx}

\vfill
\begin{center}
Updated \cvupdated
\end{center}

\end{document}
